\scanpage{3}
soient $\underline{\unsure{S|\mathcal S}}/S$ et $\underline{\unsure{\mathcal R|\mathscr R|R}}/S$ les catégories des préschémas (\unsure{resp$^{\mathrm{t}}$|respt.} réseaux) sur $S$, $\underline{S}_0/S$ la \unsure{ss-catégorie|ss. catégorie} de $\underline{S}/S$ formée des \inserted{$S$-}préschémas loc. de type fini. Alors \deleted{\unsure{$\rho$ :|$p$}}
\[
\unsure{\rho|P|p}:\underline{S}_0/S\longrightarrow\underline{\mathcal R}/S
\]
est \uline{pleinement fidèle}\unsure{,|.}
\deleted{En effet, un hom \deleted{\unsure{appl|eff}} de réseaux $\rho(X)\longrightarrow\rho(Y)$ est continu pour les topologies sous-jacentes, car on est ramené au cas où $X$ est affine donc noethérien, \unsure{donc|alors} les pts de $X$ sont spécialisations d'un \unsure{nb|nbre|ens.} fini d'entre eux, donc les pts de $=$ les pts de $f(X)$ \unsure{- - - -}.}
% slanted
\margin{est-il inutile de mettre des conditions de finitude sur $\unsure{Y|S}$ ??}

\deleted{Plus généralement, si $Y$ est \inserted{\unsure{loc|loc.}} de type fini sur $S$, \deleted{\unsure{des préschémas|deux préschémas}} \inserted{$X$ quelconque}, alors}
\[
\deleted{\operatorname{Hom}_{S}(X,Y)\xrightarrow{\sim}H_{S}(\rho(X),\rho(Y))}
\]
\deleted{[bijectif]. En effet, \inserted{$f\in\operatorname{Hom}_{S}\unsure{(|}\rho(X),\rho(Y))$. (\unsure{évidemt|d'ailleurs} \unsure{tjs|tjrs} \uline{inj}.} soit $x\in X$, $y=f(x)$, considérons $\underline{O}_y\to\underline{O}_x$, il définit un hom d'un voisinage \inserted{\deleted{\unsure{ouvert}}} $U$ de $x$ dans un voisinage \unsure{ouvert} $V$ de $y$.}
% archive sheet number: 2
